% Options for packages loaded elsewhere
% Options for packages loaded elsewhere
\PassOptionsToPackage{unicode}{hyperref}
\PassOptionsToPackage{hyphens}{url}
\PassOptionsToPackage{dvipsnames,svgnames,x11names}{xcolor}
%
\documentclass[
  letterpaper,
  DIV=11,
  numbers=noendperiod]{scrartcl}
\usepackage{xcolor}
\usepackage{amsmath,amssymb}
\setcounter{secnumdepth}{-\maxdimen} % remove section numbering
\usepackage{iftex}
\ifPDFTeX
  \usepackage[T1]{fontenc}
  \usepackage[utf8]{inputenc}
  \usepackage{textcomp} % provide euro and other symbols
\else % if luatex or xetex
  \usepackage{unicode-math} % this also loads fontspec
  \defaultfontfeatures{Scale=MatchLowercase}
  \defaultfontfeatures[\rmfamily]{Ligatures=TeX,Scale=1}
\fi
\usepackage{lmodern}
\ifPDFTeX\else
  % xetex/luatex font selection
\fi
% Use upquote if available, for straight quotes in verbatim environments
\IfFileExists{upquote.sty}{\usepackage{upquote}}{}
\IfFileExists{microtype.sty}{% use microtype if available
  \usepackage[]{microtype}
  \UseMicrotypeSet[protrusion]{basicmath} % disable protrusion for tt fonts
}{}
\makeatletter
\@ifundefined{KOMAClassName}{% if non-KOMA class
  \IfFileExists{parskip.sty}{%
    \usepackage{parskip}
  }{% else
    \setlength{\parindent}{0pt}
    \setlength{\parskip}{6pt plus 2pt minus 1pt}}
}{% if KOMA class
  \KOMAoptions{parskip=half}}
\makeatother
% Make \paragraph and \subparagraph free-standing
\makeatletter
\ifx\paragraph\undefined\else
  \let\oldparagraph\paragraph
  \renewcommand{\paragraph}{
    \@ifstar
      \xxxParagraphStar
      \xxxParagraphNoStar
  }
  \newcommand{\xxxParagraphStar}[1]{\oldparagraph*{#1}\mbox{}}
  \newcommand{\xxxParagraphNoStar}[1]{\oldparagraph{#1}\mbox{}}
\fi
\ifx\subparagraph\undefined\else
  \let\oldsubparagraph\subparagraph
  \renewcommand{\subparagraph}{
    \@ifstar
      \xxxSubParagraphStar
      \xxxSubParagraphNoStar
  }
  \newcommand{\xxxSubParagraphStar}[1]{\oldsubparagraph*{#1}\mbox{}}
  \newcommand{\xxxSubParagraphNoStar}[1]{\oldsubparagraph{#1}\mbox{}}
\fi
\makeatother

\usepackage{color}
\usepackage{fancyvrb}
\newcommand{\VerbBar}{|}
\newcommand{\VERB}{\Verb[commandchars=\\\{\}]}
\DefineVerbatimEnvironment{Highlighting}{Verbatim}{commandchars=\\\{\}}
% Add ',fontsize=\small' for more characters per line
\usepackage{framed}
\definecolor{shadecolor}{RGB}{241,243,245}
\newenvironment{Shaded}{\begin{snugshade}}{\end{snugshade}}
\newcommand{\AlertTok}[1]{\textcolor[rgb]{0.68,0.00,0.00}{#1}}
\newcommand{\AnnotationTok}[1]{\textcolor[rgb]{0.37,0.37,0.37}{#1}}
\newcommand{\AttributeTok}[1]{\textcolor[rgb]{0.40,0.45,0.13}{#1}}
\newcommand{\BaseNTok}[1]{\textcolor[rgb]{0.68,0.00,0.00}{#1}}
\newcommand{\BuiltInTok}[1]{\textcolor[rgb]{0.00,0.23,0.31}{#1}}
\newcommand{\CharTok}[1]{\textcolor[rgb]{0.13,0.47,0.30}{#1}}
\newcommand{\CommentTok}[1]{\textcolor[rgb]{0.37,0.37,0.37}{#1}}
\newcommand{\CommentVarTok}[1]{\textcolor[rgb]{0.37,0.37,0.37}{\textit{#1}}}
\newcommand{\ConstantTok}[1]{\textcolor[rgb]{0.56,0.35,0.01}{#1}}
\newcommand{\ControlFlowTok}[1]{\textcolor[rgb]{0.00,0.23,0.31}{\textbf{#1}}}
\newcommand{\DataTypeTok}[1]{\textcolor[rgb]{0.68,0.00,0.00}{#1}}
\newcommand{\DecValTok}[1]{\textcolor[rgb]{0.68,0.00,0.00}{#1}}
\newcommand{\DocumentationTok}[1]{\textcolor[rgb]{0.37,0.37,0.37}{\textit{#1}}}
\newcommand{\ErrorTok}[1]{\textcolor[rgb]{0.68,0.00,0.00}{#1}}
\newcommand{\ExtensionTok}[1]{\textcolor[rgb]{0.00,0.23,0.31}{#1}}
\newcommand{\FloatTok}[1]{\textcolor[rgb]{0.68,0.00,0.00}{#1}}
\newcommand{\FunctionTok}[1]{\textcolor[rgb]{0.28,0.35,0.67}{#1}}
\newcommand{\ImportTok}[1]{\textcolor[rgb]{0.00,0.46,0.62}{#1}}
\newcommand{\InformationTok}[1]{\textcolor[rgb]{0.37,0.37,0.37}{#1}}
\newcommand{\KeywordTok}[1]{\textcolor[rgb]{0.00,0.23,0.31}{\textbf{#1}}}
\newcommand{\NormalTok}[1]{\textcolor[rgb]{0.00,0.23,0.31}{#1}}
\newcommand{\OperatorTok}[1]{\textcolor[rgb]{0.37,0.37,0.37}{#1}}
\newcommand{\OtherTok}[1]{\textcolor[rgb]{0.00,0.23,0.31}{#1}}
\newcommand{\PreprocessorTok}[1]{\textcolor[rgb]{0.68,0.00,0.00}{#1}}
\newcommand{\RegionMarkerTok}[1]{\textcolor[rgb]{0.00,0.23,0.31}{#1}}
\newcommand{\SpecialCharTok}[1]{\textcolor[rgb]{0.37,0.37,0.37}{#1}}
\newcommand{\SpecialStringTok}[1]{\textcolor[rgb]{0.13,0.47,0.30}{#1}}
\newcommand{\StringTok}[1]{\textcolor[rgb]{0.13,0.47,0.30}{#1}}
\newcommand{\VariableTok}[1]{\textcolor[rgb]{0.07,0.07,0.07}{#1}}
\newcommand{\VerbatimStringTok}[1]{\textcolor[rgb]{0.13,0.47,0.30}{#1}}
\newcommand{\WarningTok}[1]{\textcolor[rgb]{0.37,0.37,0.37}{\textit{#1}}}

\usepackage{longtable,booktabs,array}
\newcounter{none} % for unnumbered tables
\usepackage{calc} % for calculating minipage widths
% Correct order of tables after \paragraph or \subparagraph
\usepackage{etoolbox}
\makeatletter
\patchcmd\longtable{\par}{\if@noskipsec\mbox{}\fi\par}{}{}
\makeatother
% Allow footnotes in longtable head/foot
\IfFileExists{footnotehyper.sty}{\usepackage{footnotehyper}}{\usepackage{footnote}}
\makesavenoteenv{longtable}
\usepackage{graphicx}
\makeatletter
\newsavebox\pandoc@box
\newcommand*\pandocbounded[1]{% scales image to fit in text height/width
  \sbox\pandoc@box{#1}%
  \Gscale@div\@tempa{\textheight}{\dimexpr\ht\pandoc@box+\dp\pandoc@box\relax}%
  \Gscale@div\@tempb{\linewidth}{\wd\pandoc@box}%
  \ifdim\@tempb\p@<\@tempa\p@\let\@tempa\@tempb\fi% select the smaller of both
  \ifdim\@tempa\p@<\p@\scalebox{\@tempa}{\usebox\pandoc@box}%
  \else\usebox{\pandoc@box}%
  \fi%
}
% Set default figure placement to htbp
\def\fps@figure{htbp}
\makeatother





\setlength{\emergencystretch}{3em} % prevent overfull lines

\providecommand{\tightlist}{%
  \setlength{\itemsep}{0pt}\setlength{\parskip}{0pt}}



 


\KOMAoption{captions}{tableheading}
\makeatletter
\@ifpackageloaded{caption}{}{\usepackage{caption}}
\AtBeginDocument{%
\ifdefined\contentsname
  \renewcommand*\contentsname{Table of contents}
\else
  \newcommand\contentsname{Table of contents}
\fi
\ifdefined\listfigurename
  \renewcommand*\listfigurename{List of Figures}
\else
  \newcommand\listfigurename{List of Figures}
\fi
\ifdefined\listtablename
  \renewcommand*\listtablename{List of Tables}
\else
  \newcommand\listtablename{List of Tables}
\fi
\ifdefined\figurename
  \renewcommand*\figurename{Figure}
\else
  \newcommand\figurename{Figure}
\fi
\ifdefined\tablename
  \renewcommand*\tablename{Table}
\else
  \newcommand\tablename{Table}
\fi
}
\@ifpackageloaded{float}{}{\usepackage{float}}
\floatstyle{ruled}
\@ifundefined{c@chapter}{\newfloat{codelisting}{h}{lop}}{\newfloat{codelisting}{h}{lop}[chapter]}
\floatname{codelisting}{Listing}
\newcommand*\listoflistings{\listof{codelisting}{List of Listings}}
\makeatother
\makeatletter
\makeatother
\makeatletter
\@ifpackageloaded{caption}{}{\usepackage{caption}}
\@ifpackageloaded{subcaption}{}{\usepackage{subcaption}}
\makeatother
\usepackage{bookmark}
\IfFileExists{xurl.sty}{\usepackage{xurl}}{} % add URL line breaks if available
\urlstyle{same}
\hypersetup{
  pdftitle={Lab: Paper Helicopters},
  colorlinks=true,
  linkcolor={blue},
  filecolor={Maroon},
  citecolor={Blue},
  urlcolor={Blue},
  pdfcreator={LaTeX via pandoc}}


\title{Lab: Paper Helicopters}
\usepackage{etoolbox}
\makeatletter
\providecommand{\subtitle}[1]{% add subtitle to \maketitle
  \apptocmd{\@title}{\par {\large #1 \par}}{}{}
}
\makeatother
\subtitle{Factorial Design and ANOVA}
\author{}
\date{}
\begin{document}
\maketitle


\subsection{Introduction}\label{introduction}

The goal of this exercise is to get a better grip on the concept of
Factorial Designs, their advantages, and their corresponding tests and
statistical analyses. We study this in the context of an engineering
experiment in aerodynamics, before proceeding to conduct our versions of
such an experiment using paper helicopters! Finally, we implement the
ANOVA (Analysis of Variance) framework from lecture for our collected
data.\\

\subsection{The Experiment}\label{the-experiment}

We first look at a recent real-world application of factorial designs in
aerodynamics experiments from 2020, which investigates the effect of
surface roughness and other environmental factors on airfoil
performance. Read through the paper, trying to identify the details of
the design such as the factors and levels, and the motivation for a full
factorial design. Also try to interpret the results and understand their
significance.\\

Answer the following questions in full sentences.

\begin{enumerate}
\def\labelenumi{\arabic{enumi}.}
\tightlist
\item
  What is the main motivation for conducting a full factorial design in
  such a setting? What would changing one factor at a time fail to
  detect in this study?
\end{enumerate}

\hfill\break
\hfill\break
\hfill\break
\hfill\break
\hfill\break
\hfill\break
\hfill\break
\hfill\break

\begin{enumerate}
\def\labelenumi{\arabic{enumi}.}
\setcounter{enumi}{1}
\tightlist
\item
  Which of the factors are significant from the ANOVA analysis?
  Interpret what it means if two factors and their interaction are all
  significant? Are they independent?
\end{enumerate}

\hfill\break
\hfill\break
\hfill\break
\hfill\break
\hfill\break
\hfill\break
\hfill\break
\hfill\break

\subsection{Conducting Your Own
Experiment}\label{conducting-your-own-experiment}

Conduct your own factorial design experiment now with paper helicopters!
Fold the provided templates after adjusting for different wing heights
and helicopter weights (via paper clips) to analyze whether they have a
significant effect on the duration of flight. Choose the levels that you
wish to study in your experiment, and feel free to add other factors to
your design.

Keep the following questions in mind, and write answers explaining your
choice in each.

\begin{enumerate}
\def\labelenumi{\arabic{enumi}.}
\setcounter{enumi}{2}
\tightlist
\item
  List out the factors and levels that you plan on including in your
  design experiment. Guess whether these factors would increase or
  decrease the flight time.
\end{enumerate}

\hfill\break
\hfill\break
\hfill\break
\hfill\break
\hfill\break
\hfill\break

\begin{enumerate}
\def\labelenumi{\arabic{enumi}.}
\setcounter{enumi}{3}
\tightlist
\item
  How many runs or replicates are you planning to conduct for each
  combination of factors? Will different people drop the helicopter or
  time the experiment for each run? Would this be a factor or would it
  simply be different replicates?
\end{enumerate}

\hfill\break
\hfill\break
\hfill\break
\hfill\break
\hfill\break
\hfill\break
\hfill\break

\begin{enumerate}
\def\labelenumi{\arabic{enumi}.}
\setcounter{enumi}{4}
\tightlist
\item
  Is the same person going to fold every helicopter? What would happen
  if you drew out the designs instead of using a template? Does it make
  your experiment more robust? What does it mean for power?
\end{enumerate}

\hfill\break
\hfill\break
\hfill\break
\hfill\break
\hfill\break
\hfill\break
\hfill\break
\hfill\break

\begin{enumerate}
\def\labelenumi{\arabic{enumi}.}
\setcounter{enumi}{5}
\tightlist
\item
  Where are you going to drop the helicopter from? Suppose your
  engineering experiment involved very expensive equipment (such as
  rockets or drones) and a large cost associated with each run. How
  would you decide the drop height?
\end{enumerate}

\hfill\break
\hfill\break
\hfill\break
\hfill\break
\hfill\break
\hfill\break
\hfill\break
\hfill\break

\begin{enumerate}
\def\labelenumi{\arabic{enumi}.}
\setcounter{enumi}{6}
\tightlist
\item
  Sketch the dataframe that will guide your experiment. Include the
  following columns: an ID to keep track of each measurement unit, the
  index of the replicate, each experimental factor that will be used and
  its level, an index to track the treatment group number (\texttt{d}),
  and the measured response.
\end{enumerate}

\hfill\break
\hfill\break
\hfill\break
\hfill\break
\hfill\break
\hfill\break
\hfill\break
\hfill\break
\hfill\break
\hfill\break
\hfill\break
\hfill\break
\hfill\break
\hfill\break
\hfill\break
\hfill\break
\hfill\break
\hfill\break
\hfill\break
\hfill\break
\hfill\break
\hfill\break
\hfill\break
\hfill\break

\begin{enumerate}
\def\labelenumi{\arabic{enumi}.}
\setcounter{enumi}{7}
\tightlist
\item
  Should the runs be conducted in the order you listed in your design
  matrix? What might happen if you run all ``heavy'' helicopters first
  and all ``light'' helicopters later?
\end{enumerate}

\hfill\break
\hfill\break
\hfill\break
\hfill\break
\hfill\break
\hfill\break
\hfill\break
\hfill\break

\subsection{Data Analysis}\label{data-analysis}

We will now analyze our own collected data. Ensure that you tabulate
your data into a \texttt{.csv} file. This can be done by entering your
data into Google Sheets or Microsoft Excel before saving as
\texttt{.csv}.

As earlier, use the \texttt{readr} package, which is included inside the
\texttt{tidyverse}. If you haven't installed the tidyverse before, you
can do so by running \texttt{install.packages("tidyverse")} once.

\begin{Shaded}
\begin{Highlighting}[]
\CommentTok{\# load tidyverse (includes readr for CSVs)}
\FunctionTok{library}\NormalTok{(tidyverse)}
\end{Highlighting}
\end{Shaded}

\begin{Shaded}
\begin{Highlighting}[]
\CommentTok{\# load your collected data into R}

\CommentTok{\#my\_data \textless{}{-} }
\end{Highlighting}
\end{Shaded}

\subsubsection{EDA and Scatterplots}\label{eda-and-scatterplots}

\begin{enumerate}
\def\labelenumi{\arabic{enumi}.}
\setcounter{enumi}{8}
\tightlist
\item
  Explore how your helicopters' flight times differed by wing width.
  Create a beehive plot using the \texttt{ggplot2} and
  \texttt{ggbeeswarm} packages in \texttt{R}. If you don't have these
  installed, install them by running
  \texttt{install.packages("ggplot2")} and
  \texttt{install.packages("ggbeeswarm")} once.
\end{enumerate}

Stick to a single level of weight while making this for height, or add
levels of weight by color coding them!

\begin{Shaded}
\begin{Highlighting}[]
\CommentTok{\# Load necessary packages}

\FunctionTok{library}\NormalTok{(ggplot2)}
\FunctionTok{library}\NormalTok{(ggbeeswarm)}

\CommentTok{\# Beeswarm Plot (Fill in)}

\CommentTok{\# my\_data |\textgreater{}}
\CommentTok{\#   ggplot(aes(x = , y = )) +}
\CommentTok{\#   geom\_beeswarm(cex = 3, size = 3)}
\end{Highlighting}
\end{Shaded}

Next, calculate the group means. You may choose to this for any one
factor or group by both. This requires the \texttt{dplyr} package, which
you can install using \texttt{install.packages("dplyr")}. Also calculate
the difference of each group mean with the overall mean. You may plot
this to see if it is meaningful. Is it linear, non-linear etc?

\begin{Shaded}
\begin{Highlighting}[]
\DocumentationTok{\#\#\# Load necessary packages}

\FunctionTok{library}\NormalTok{(dplyr)}

\DocumentationTok{\#\#\# Calculate Group Means}

\CommentTok{\# group\_means \textless{}{-} my\_data |\textgreater{}}
\CommentTok{\#   group\_by() |\textgreater{}}
\CommentTok{\#   summarize(avg = mean())}
\CommentTok{\# group\_means}
\end{Highlighting}
\end{Shaded}

\begin{Shaded}
\begin{Highlighting}[]
\DocumentationTok{\#\#\# Group Means {-} Overall Means}

\CommentTok{\# overall\_mean \textless{}{-} mean()}

\CommentTok{\# effects\_df \textless{}{-} group\_means |\textgreater{}}
\CommentTok{\#   mutate(effect = avg {-} overall\_mean)}

\CommentTok{\# effects\_df$effect}
\end{Highlighting}
\end{Shaded}

\begin{Shaded}
\begin{Highlighting}[]
\DocumentationTok{\#\#\# Plot:}

\CommentTok{\# ggplot(effects\_df,}
\CommentTok{\#        aes(x = width,}
\CommentTok{\#            y = effect,}
\CommentTok{\#            group = 1)) +}
\CommentTok{\#   geom\_point(size = 3) +}
\CommentTok{\#   geom\_line(linewidth = 1) +}
\CommentTok{\#   geom\_hline(yintercept = 0, linetype = "dashed") +}
\CommentTok{\#   labs(x = "Wing Width",}
\CommentTok{\#        y = "Difference from Overall Mean")}
\end{Highlighting}
\end{Shaded}

\subsubsection{Randomization Based
Inference}\label{randomization-based-inference}

\begin{enumerate}
\def\labelenumi{\arabic{enumi}.}
\setcounter{enumi}{9}
\tightlist
\item
  Next, conduct a randomization based test for the F-statistic to check
  whether the width of the wings turned out to be significant in
  determining flight times. Again, note that you will have to do this
  after fixing a level for weight.
\end{enumerate}

For reference, see the
\href{https://stat158.berkeley.edu/spring-2026/13-f-inference/slides.html\#/title-slide}{Stat
158 F-inference slides page}.

\begin{Shaded}
\begin{Highlighting}[]
\DocumentationTok{\#\#\# Selecting the weight factor}

\CommentTok{\# my\_data\_subset \textless{}{-} my\_data |\textgreater{}}
\CommentTok{\#   filter(weight == "light")}
\end{Highlighting}
\end{Shaded}

\begin{Shaded}
\begin{Highlighting}[]
\DocumentationTok{\#\#\# Function to simulate under the Null}

\NormalTok{shuffle4\_f }\OtherTok{\textless{}{-}} \ControlFlowTok{function}\NormalTok{(y, z) \{}
\NormalTok{  Z }\OtherTok{\textless{}{-}} \FunctionTok{sample}\NormalTok{(z) }\CommentTok{\# source of randomness}

\NormalTok{  my\_data\_subset }\OtherTok{\textless{}{-}} \FunctionTok{data.frame}\NormalTok{(}\AttributeTok{flight\_time =}\NormalTok{ y,}
                           \AttributeTok{width =}\NormalTok{ Z)}
  
\NormalTok{  avg\_var }\OtherTok{\textless{}{-}}\NormalTok{ my\_data\_subset }\SpecialCharTok{|\textgreater{}}
    \FunctionTok{group\_by}\NormalTok{(width) }\SpecialCharTok{|\textgreater{}}
    \FunctionTok{summarize}\NormalTok{(}\AttributeTok{var =} \FunctionTok{var}\NormalTok{(flight\_time)) }\SpecialCharTok{|\textgreater{}}
    \FunctionTok{summarize}\NormalTok{(}\AttributeTok{avg\_var =} \FunctionTok{mean}\NormalTok{(var)) }\SpecialCharTok{|\textgreater{}}
    \FunctionTok{pull}\NormalTok{(avg\_var)}
\NormalTok{  var\_avg }\OtherTok{\textless{}{-}}\NormalTok{ my\_data\_subset }\SpecialCharTok{|\textgreater{}}
    \FunctionTok{group\_by}\NormalTok{(width) }\SpecialCharTok{|\textgreater{}}
    \FunctionTok{summarize}\NormalTok{(}\AttributeTok{avg =} \FunctionTok{mean}\NormalTok{(flight\_time)) }\SpecialCharTok{|\textgreater{}}
    \FunctionTok{summarize}\NormalTok{(}\AttributeTok{var\_avg =} \FunctionTok{var}\NormalTok{(avg)) }\SpecialCharTok{|\textgreater{}}
    \FunctionTok{pull}\NormalTok{(var\_avg)}

\NormalTok{  var\_avg }\SpecialCharTok{/}\NormalTok{ avg\_var}
\NormalTok{\}}
\end{Highlighting}
\end{Shaded}

\begin{Shaded}
\begin{Highlighting}[]
\DocumentationTok{\#\#\# Calculating the observed value of F {-} statistic}

\NormalTok{avg\_var }\OtherTok{\textless{}{-}}\NormalTok{ my\_data\_subset }\SpecialCharTok{|\textgreater{}}
  \FunctionTok{group\_by}\NormalTok{(width) }\SpecialCharTok{|\textgreater{}}
  \FunctionTok{summarize}\NormalTok{(}\AttributeTok{var =} \FunctionTok{var}\NormalTok{(flight\_time)) }\SpecialCharTok{|\textgreater{}}
  \FunctionTok{summarize}\NormalTok{(}\AttributeTok{avg\_var =} \FunctionTok{mean}\NormalTok{(var)) }\SpecialCharTok{|\textgreater{}}
  \FunctionTok{pull}\NormalTok{(avg\_var)}

\NormalTok{var\_avg }\OtherTok{\textless{}{-}}\NormalTok{ my\_data\_subset }\SpecialCharTok{|\textgreater{}}
  \FunctionTok{group\_by}\NormalTok{(width) }\SpecialCharTok{|\textgreater{}}
  \FunctionTok{summarize}\NormalTok{(}\AttributeTok{avg =} \FunctionTok{mean}\NormalTok{(flight\_time)) }\SpecialCharTok{|\textgreater{}}
  \FunctionTok{summarize}\NormalTok{(}\AttributeTok{var\_avg =} \FunctionTok{var}\NormalTok{(avg)) }\SpecialCharTok{|\textgreater{}}
  \FunctionTok{pull}\NormalTok{(var\_avg)}

\NormalTok{f\_stat }\OtherTok{\textless{}{-}}\NormalTok{ var\_avg }\SpecialCharTok{/}\NormalTok{ avg\_var}
\end{Highlighting}
\end{Shaded}

\begin{Shaded}
\begin{Highlighting}[]
\DocumentationTok{\#\#\# Plotting the Null Distribution}

\CommentTok{\# null\_stats \textless{}{-} replicate(500, shuffle4\_f(y = , z = ))}
\NormalTok{null }\OtherTok{\textless{}{-}} \FunctionTok{data.frame}\NormalTok{(}\AttributeTok{stats =}\NormalTok{ null\_stats)}

\FunctionTok{ggplot}\NormalTok{(null, }\FunctionTok{aes}\NormalTok{(}\AttributeTok{x =}\NormalTok{ stats)) }\SpecialCharTok{+}
  \FunctionTok{geom\_histogram}\NormalTok{() }\SpecialCharTok{+}
  \FunctionTok{geom\_vline}\NormalTok{(}\AttributeTok{xintercept =}\NormalTok{ f\_stat, }\AttributeTok{color =} \StringTok{"tomato"}\NormalTok{, }\AttributeTok{lwd =} \DecValTok{2}\NormalTok{)}

\DocumentationTok{\#\#\# Calculating p{-}value}

\NormalTok{null }\SpecialCharTok{|\textgreater{}}
  \FunctionTok{summarize}\NormalTok{(}\AttributeTok{pval =} \FunctionTok{mean}\NormalTok{(stats }\SpecialCharTok{\textgreater{}}\NormalTok{ f\_stat))}
\end{Highlighting}
\end{Shaded}

\subsubsection{Model Based Inference}\label{model-based-inference}

\begin{enumerate}
\def\labelenumi{\arabic{enumi}.}
\setcounter{enumi}{10}
\tightlist
\item
  We now move into two factor results using the Model Based approach.
\end{enumerate}

What are the assumptions of the Model-based F-test used here?

\hfill\break
\hfill\break
\hfill\break
\hfill\break

Now, proceed to calculate the ANOVA Table as discussed in Lecture. Note
that you can use \texttt{+} to add more than 1 factor, and \texttt{*} to
specify interactions such as
\texttt{aov(y\ \textasciitilde{}\ x1\ +\ x2\ +\ x1:x2)}.

\begin{Shaded}
\begin{Highlighting}[]
\DocumentationTok{\#\#\# Calculating the ANOVA table}

\CommentTok{\# helicopter\_anova \textless{}{-} aov(, data = my\_data)}
\CommentTok{\# anova\_table \textless{}{-} summary(helicopter\_anova)}
\CommentTok{\# anova\_table}
\end{Highlighting}
\end{Shaded}

What factors do you find significant? Do they align with the EDA and the
Randomization based tests?

\hfill\break
\hfill\break
\hfill\break
\hfill\break

\subsubsection{Interaction Plots}\label{interaction-plots}

\begin{enumerate}
\def\labelenumi{\arabic{enumi}.}
\setcounter{enumi}{11}
\tightlist
\item
  Finally, we plot another simple diagnostic to check the pattern of any
  interaction between the two factors, wing width and helicopter weight
  (paper clips). We do this by plotting separate line graphs of flight
  time by width for both ``light'' and ``heavy'' helicopters.
\end{enumerate}

If the lines are roughly parallel, this indicates that the effect of
``heaviness'' doesn't change with wing width. However if they aren't the
plot gives you an idea of how this effect of weight changes with width,
i.e it gives you an idea of how the interaction effect behaves.

\begin{Shaded}
\begin{Highlighting}[]
\DocumentationTok{\#\#\# First compute group means for each combination}
\DocumentationTok{\#\#\# Fill in appropriate column names}

\NormalTok{interaction\_means }\OtherTok{\textless{}{-}}\NormalTok{ my\_data }\SpecialCharTok{|\textgreater{}}
  \FunctionTok{group\_by}\NormalTok{(width, weight) }\SpecialCharTok{|\textgreater{}}
  \FunctionTok{summarize}\NormalTok{(}\AttributeTok{avg\_flight\_time =} \FunctionTok{mean}\NormalTok{(flight\_time))}

\NormalTok{interaction\_means}
\end{Highlighting}
\end{Shaded}

\begin{Shaded}
\begin{Highlighting}[]
\DocumentationTok{\#\#\# Plot the interaction}
\DocumentationTok{\#\#\# Fill in appropriate column names}

\FunctionTok{ggplot}\NormalTok{(interaction\_means,}
       \FunctionTok{aes}\NormalTok{(}\AttributeTok{x =}\NormalTok{ width,}
           \AttributeTok{y =}\NormalTok{ avg\_flight\_time,}
           \AttributeTok{group =}\NormalTok{ weight,}
           \AttributeTok{color =}\NormalTok{ weight)) }\SpecialCharTok{+}
  \FunctionTok{geom\_point}\NormalTok{(}\AttributeTok{size =} \DecValTok{3}\NormalTok{) }\SpecialCharTok{+}
  \FunctionTok{geom\_line}\NormalTok{(}\AttributeTok{linewidth =} \DecValTok{1}\NormalTok{) }\SpecialCharTok{+}
  \FunctionTok{labs}\NormalTok{(}\AttributeTok{x =} \StringTok{"Wing Width"}\NormalTok{,}
       \AttributeTok{y =} \StringTok{"Average Flight Time"}\NormalTok{,}
       \AttributeTok{color =} \StringTok{"Paperclip Weight"}\NormalTok{)}
\end{Highlighting}
\end{Shaded}

Are the two lines parallel? Is this consistent with the ANOVA table
results above?

\hfill\break
\hfill\break
\hfill\break
\hfill\break

Answer in brief: how did the structure of your experiment affect what
conclusions you were allowed to draw from its final results?

\hfill\break
\hfill\break
\hfill\break
\hfill\break




\end{document}
